\documentclass[12pt]{article}

% Layout.
\usepackage[top=1in, bottom=0.75in, left=0.75in, right=0.75in, headheight=1in, headsep=6pt]{geometry}

% Fonts.
\usepackage{mathptmx}

\usepackage[scaled=0.86]{helvet}
\renewcommand{\emph}[1]{\textsf{\textbf{#1}}}

% Misc packages.
\usepackage{amsmath,amssymb,latexsym,amsthm}
\usepackage{graphicx,hyperref}
\usepackage{array}
\usepackage{xcolor}
\usepackage{multicol}
\usepackage{tabularx,colortbl}
\usepackage{enumitem}
\usepackage{soul,multicol}
\usepackage{setspace}


\hypersetup{
    colorlinks=true,
    linkcolor=blue,
    filecolor=magenta,      
    urlcolor=blue,
    pdftitle={Proofs Worksheet},
    pdfpagemode=FullScreen,
    }

\def\mailto#1{\href{mailto:#1}{#1}}

%% special math symbols
\newcommand{\N}{\mathbb{N}}
\newcommand{\Z}{\mathbb{Z}}
\newcommand{\R}{\mathbb{R}}
\newcommand{\Q}{\mathbb{Q}}
\newcommand{\sse}{\subseteq}
\newcommand{\mcp}{\mathcal{P}}
\newcommand{\ol}[1]{\overline{#1}}

%other specials
\newcommand{\be}{\begin{enumerate}}
\newcommand{\ee}{\end{enumerate}}
\newcommand{\se}{\subseteq}
\newcommand{\spe}{\supseteq}
\newcommand{\ds}{\displaystyle}
\newcommand{\lcm}{\text{lcm}}
%\newcommand{\gcd}{\text{gcd}}

% Paragraph spacing
\parindent 0pt
\parskip 6pt plus 1pt
\def\tableindent{\hskip 0.5 in}
\def\ts{\hskip 1.5 em}

%header
\usepackage{fancyhdr}
\pagestyle{fancy} 
\lhead{\large\sf\textbf{MATH 265: Introduction to Mathematical Proofs}}
\rhead{\large\sf\textbf{Worksheet 25: Ch 14 \S 3-4}}

\newcommand{\localhead}[1]{\par\smallskip\textbf{#1}\nobreak\\}%
\def\heading#1{\localhead{\large\emph{#1}}}
\def\subheading#1{\localhead{\emph{#1}}}

\newenvironment{clist}%
{\bgroup\parskip 0pt\begin{list}{$\bullet$}{\partopsep 4pt\topsep 0pt\itemsep -2pt}}%
{\end{list}\egroup}%

\begin{document}
\begin{enumerate}
\item Thm 14.7: Let $A$ be a set. Then \quad \fbox{$|A|$ \quad \quad $|\mathcal{P}(A)|.$}\\
\vfill
\item Thm 14.8: Let $A$ be a countably infinite set and $B$ be an infinite set such that $B\subseteq A.$ Then \\
\vfill
\item Thm 14.10 (The Cantor-Bernstein-Schröder Theorem)\\
Let $f:A\to B$ and $g:B \to A$ be injective. Then, 
\vfill
\item Thm 14.9: Let $B \subseteq A$ such that $B$ is uncountable. Then, \\
\vfill
\item Let $f(x)=\frac{1}{1+e^x}$
	\begin{enumerate}
	\item Suppose that the domain of $f$ is $\R.$ Determine the range of $f$ and justify your conclusion.
	\vfill
	\item Is $f(x)$ injective? Justify your conclusion.
	\vfill
	\item Now use $f(x)$ (and your work above) to prove two sets have the same cardinality.
	\vfill
	\end{enumerate}
\newpage
\item Show that $|(0,1)|=|[0,1)|$ by finding a bijection between them.\\
\vfill
\item Show that $|(0,1)|=|[0,1)|$ using the Cantor-Bernstein-Schröder Theorem.\\
\vfill
\item Thm 14.11 $|\R|=|\mathcal{P}(\N)|$\\
\vspace{5in}
\newpage
\end{enumerate}
\end{document}










